\chapter{双重不对称性：任务可辨识度与主体控制偏置}
\label{chap:dual_asymmetry}

上一章说明，带身份的状态消息如何在受约束网络中传播；本章进一步讨论传播方向为何通常不对称。同一对事实关系，正向生成与逆向诊断可能具有不同的可辨识度、样本需求和搜索复杂度；同一条可行路径，也可能因为目标、承诺、风险和资源状态不同而获得不同的控制优先级。前一种差异来自任务与观测机制，后一种差异来自主体的规范状态，两者可以相互作用，却不能被合并成同一种“阻力”。我们因此把原有的黎曼曲率、热力学箭头和认知洛伦兹力降为辅助意象，以条件化逆问题、版本化控制偏置和可审计联合代价重建双重不对称性。这样的模型既能解释生成、诊断、科研、家务与放弃等具体案例，也能说明何时需要增加证据、改变任务分解、调节控制强度或执行认知卸载。

\section{事实不对称：正向映射、逆问题与可辨识度}
\label{sec:factual_asymmetry}

我们先定义所谓事实不对称。设环境或问题状态为 $x\in\mathcal{X}$，可观测结果为 $y\in\mathcal{Y}$，任务条件为 $c\in\mathcal{C}$。正向模型由条件核 $K_c(y\mid x)$ 给出，它描述在状态 $x$ 和条件 $c$ 下观察到 $y$ 的分布；若系统是确定性的，则退化为映射 $y=F_c(x)$。正向任务是在给定 $x,c$ 时预测或生成 $y$，逆向任务是在给定 $y,c$ 时识别可能的 $x$。两者使用同一观测机制，却不自动具有相同难度。只要存在 $x_1\neq x_2$ 使 $K_c(\cdot\mid x_1)$ 与 $K_c(\cdot\mid x_2)$ 相同或高度重叠，单次观测就不能唯一恢复来源；这首先是信息不足与模型不可辨识，而不是逆向运动违反某种热力学方向。

对有限候选集，可以用区分裕量描述局部可辨识度。给定候选 $x_i$ 的预测分布 $p_i(y)=K_c(y\mid x_i)$，定义
\begin{equation}
\Delta_c(x_i)=\min_{j\neq i}D\!\left(p_i\,\|\,p_j\right),
\end{equation}
其中 $D$ 是事先声明的统计散度。若 $\Delta_c(x_i)$ 很小，候选在当前观测协议下难以区分；增加推理步数并不会凭空创造信息，必须取得新观测、改变干预或缩小候选集合。连续参数模型中，可用局部 Fisher 信息或观测映射雅可比的奇异值诊断病态方向，但这些量依赖参数化、噪声模型和采样方案。它们说明“哪些变化难以被看见”，并不等于世界本身具有固定的认知曲率。

正向任务也未必容易。生成一个满足语法的句子可能候选很多，但生成同时满足安全规范、事实约束、制造公差和后续可维护性的方案，合法集合可能极小；此时寻找任意输出仍然困难。相反，有些逆向任务具有唯一、稳定且廉价的解，例如带校验码的数据恢复或可逆程序的输入重建。因此，“生成总比理解容易”只能作为特定观测与约束下的经验倾向。我们比较任务时至少要同时记录候选空间规模、观测噪声、约束密度、允许误差、所需置信和验证成本。

更一般地，事实任务是一组带证据的搜索过程。令 $S_k\subseteq\mathcal{X}$ 为第 $k$ 步尚未排除的候选，观测或干预 $a_k$ 产生记录 $o_k$，则候选更新为
\begin{equation}
S_{k+1}=\{x\in S_k:C(x,o_k,a_k,c_k)=1\}.
\end{equation}
这里 $C$ 是证据兼容谓词，$c_k$ 允许任务条件随时间改变。逆向诊断的有效进展体现在候选集收缩、后验校准或最坏风险下降，而不是内部活动强度上升。若一次观测同时排除大量候选，逆问题可以迅速变简单；若观测彼此相关或来自同一错误来源，表面证据数量增加也不会带来相同比例的区分力。

因果方向还必须和计算方向分开。原因先于结果是一种领域主张；从结果查询原因是一种算法任务。系统可以沿因果模型正向模拟，也可以利用贝叶斯反演、约束求解、检索或实验设计逆向推断。逆向结果通常依赖先验 $p(x\mid c)$：
\begin{equation}
p(x\mid y,c)=\frac{K_c(y\mid x)p(x\mid c)}{\int_{\mathcal{X}}K_c(y\mid x')p(x'\mid c)\,\mathrm{d}x'}.
\end{equation}
该式要求分母存在且非零；后验集中也只说明所选模型与先验下的判断，不证明真实原因被唯一找到。我们因此把“事实几何”落实为可观测的候选结构、区分度和求解代价，而不由高维、负曲率或熵增意象直接推出难度结论。

工程验证采用成对任务，而不是只比较主观体验。我们为同一机制构造正向预测、逆向诊断、带新干预的逆向诊断和约束生成四组测试，保持模型、预算与评分协议可比；记录正确率、校准、候选缩减、调用次数、延迟和失败类型。若逆向任务失败，应区分不可辨识、算法搜索不足、先验偏差与接口丢失。只有这样，所谓第一重不对称才成为可以被证伪和修复的计算性质。

\section{主体不对称：规范状态、行动资格与控制投入}
\label{sec:subjective_asymmetry}

第二重不对称发生在事实可行之后。一个行动即使逻辑上可执行，也可能不符合当前目标、承诺、授权或风险边界。我们用规范状态 $n_t=(g_t,H_t,P_t,B_t)$ 表示时刻 $t$ 的目标集合、硬约束、软偏好与资源预算；用候选行动 $a\in\mathcal{A}(x_t)$ 表示在事实状态 $x_t$ 下可考虑的操作。首先筛除违反硬约束的行动，得到资格集合
\begin{equation}
\mathcal{A}^{\mathrm{ok}}_t=\{a\in\mathcal{A}(x_t):h_j(x_t,a,n_t)\leq 0,\ j=1,\ldots,m\}.
\end{equation}
只有在这个集合内，系统才比较收益、风险、延迟和控制成本。偏好不能通过一个“场力”把越权行动重新吹回可行域，探索温度也不能给非法行动分配非零概率。

对合格行动，我们定义可审计评分，而不把欲望、活动或实际能耗混称为价值能量。设 $U_t(a)$ 是任务收益，$R_t(a)$ 是预测风险，$L_t(a)$ 是延迟，$C_t(a)$ 是计算或外部调用成本，$K_t(a)$ 是维持执行所需的控制投入。若这些量已通过声明的尺度映射为可比较的无量纲分数，可写
\begin{equation}
J_t(a)=-w_U U_t(a)+w_R R_t(a)+w_L L_t(a)+w_C C_t(a)+w_K K_t(a).
\end{equation}
若不能可靠换算，就使用词典序、Pareto 前沿或约束优化，而不是任意相减。权重属于规范状态的一部分，必须带版本、来源和适用范围；目标变化时，旧权重不能被悄悄沿用。

原章“顺流”和“逆流”的经验差异，可以由控制投入 $K_t$ 具体化。习惯动作、熟练工具和高度一致的子目标通常减少重规划、冲突消解与工作记忆占用，因此维持成本较低；戒瘾、承认错误或执行厌恶任务则可能要求持续抑制高概率候选、调用外部承诺和抵抗短期奖励，因此控制成本较高。这是关于策略竞争和资源占用的模型，不是主观价值场对路径做物理功。个体报告的“想做”或“抗拒”可作为观测之一，但还需结合中断率、重试、延迟、错误与控制动作，避免把语言自述当成唯一事实。

主体偏置并非固定人格常数。疲劳、截止时间、新证据、社会承诺和环境风险都可能改变 $n_t$。因此完整导数或离散更新必须包含显式时间项和状态依赖；在离散实现中，我们记录
\begin{equation}
n_{k+1}=\mathcal{U}_n(n_k,e_k,f_k),
\end{equation}
其中 $e_k$ 是新事件，$f_k$ 是反馈，$\mathcal{U}_n$ 必须受身份、权限和审计规则约束。一次即时偏好不应自动覆盖长期承诺；反之，长期目标也不能无限压制身体或系统安全信号。规范更新本身是一个需要证据与责任归属的过程。

主体不对称还会被结构学习放大。某条路径反复获得奖励后，检索延迟降低、默认优先级升高，未来选择看起来更“自然”；但熟悉不等于正确。系统应把路径频率、成功率、环境版本与失败条件分开保存，并在分布变化时降低旧策略资格。对于人类案例，打游戏、看爽文、家务、填表、科研和竞技都不能仅凭类别预判：同一活动对不同主体、时期与目标可能落在完全不同的控制区间。

在 AgentOS 实例中，$\Psi$ 与 SPC 提供面向当前自我和任务的压缩读出，TCE 调节合格候选内的探索与抑制，CCM 监测 $h(t)$ 的冲突、负载和重规划，Boundary 执行权限检查，Offloading 将高维护成本步骤外化到工具、文件或其他 Agent。$E$ 保存行动与反馈事件，$\Theta$ 保存较稳定的策略和承诺结构。这些机制说明规范偏置如何工程化，但 HSF-HD 只要求存在等价的目标记录、资格检查、控制接口和验证闭环，不把 AgentOS 组件规定为所有智能系统的唯一实现。

验证第二重不对称时，我们让事实难度近似相同而改变目标、授权、奖励期限和资源状态，观察策略是否按声明规则变化；再保持规范状态不变而改变任务难度，检验系统能否区分事实障碍与动机障碍。若增加证据能够解决的问题被错误归为“不想做”，或调整奖励即可解决的问题被错误归为“不会做”，诊断就失败。主体不对称的价值正在于提供可分解的修复方向，而不是为一切迟滞贴上心理摩擦的名称。

\section{联合代价与四类运行区间}
\label{sec:quadrants_of_thought}

两重不对称会在实际任务中耦合。我们令 $D_t^{\mathrm{fact}}$ 表示事实侧难度诊断向量，包括候选歧义、观测噪声、约束密度、搜索分支和验证成本；令 $D_t^{\mathrm{subj}}$ 表示主体侧诊断向量，包括目标冲突、抑制需求、承诺切换、资源紧张和风险压力。它们不是天然同单位的标量。只有在明确的评估协议 $\phi_f,\phi_s$ 下，才能得到归一化指标 $d_f=\phi_f(D_t^{\mathrm{fact}})\in[0,1]$ 与 $d_s=\phi_s(D_t^{\mathrm{subj}})\in[0,1]$，并据阈值形成四类运行区间。阈值应由任务族、系统容量和误报代价校准，不是普遍常数。

联合决策可写成受约束的多目标问题：
\begin{equation}
\min_{\pi\in\Pi_t^{\mathrm{ok}}}
\bigl(C_{\mathrm{infer}}(\pi),C_{\mathrm{control}}(\pi),R(\pi),L(\pi)\bigr),
\end{equation}
其中 $\Pi_t^{\mathrm{ok}}$ 是满足事实前置条件、权限与安全约束的策略集合。若系统采用加权标量化，必须声明量纲转换和权重；若采用词典序，则应说明哪些约束优先。目标函数下降只表示内部评分改善，不自动等于事实正确、主体满意或外部任务成功，最终仍需由环境反馈和保留任务验证。

低事实难度、低主体阻力对应“顺行区”。典型现象是步骤熟悉、反馈快速、目标一致，例如熟练操作、娱乐或自动化反应。系统不应把它称为零阻力或超导，因为仍有计算、注意、时间和真实能耗；更准确的描述是当前接口下推理成本与控制投入均低。这个区间适合保持较低调度开销，但仍需防止习惯路径在环境改变后继续自动执行。

高事实难度、低主体阻力对应“攻坚区”。科研、解谜和竞技可能落入此区：主体愿意投入，但候选歧义、长程依赖或验证成本较高。有效策略是增加证据、分解任务、建立外部工作台和阶段性检查点，而不是仅提高兴奋或探索温度。若挑战与技能匹配、反馈及时且负荷可控，系统可能进入心流；但心流是运行诊断，不证明问题已经解决，也不意味着资源消耗消失。

低事实难度、高主体阻力对应“拖滞区”。家务、填表和重复维护常被作为例子，但具体归类依主体与语境而变。此时继续增加推理能力通常收益有限，更有效的是减少启动步骤、重构奖励、明确承诺、自动化重复操作或把任务卸载到外部工具。若硬约束要求必须执行，控制器需要限制持续抑制的代价并安排恢复；不能把短期合规当作长期稳定。

高事实难度、高主体阻力对应“失稳区”。厌学、习得性无助或复杂项目崩溃可能呈现为反复失败、目标冲突、负荷过高和逃避。系统不应由一次拒绝推断永久无能力，也不应把继续加压当作唯一修复。可行措施包括降低任务分辨率、补充观测、修改环境、重议目标、引入协作者和设置安全退出。若资格集合为空，应报告不可行及冲突来源，而不是用更强控制掩盖矛盾。

四区不是人格标签，而是随时间变化的诊断状态。设进入高区阈值为 $\tau^+$，退出阈值为更低的 $\tau^-$，并要求状态持续若干窗口，可形成滞回以避免在阈值附近频繁切换。任务分解可能同时降低 $d_f$ 与 $d_s$；新风险信息也可能同时提高两者。我们因此记录区间转移的触发证据、控制动作和恢复时间，而不把四象限解释为固定拓扑相或热力学熔断。

四区诊断还应接受基线校正。同样的延迟对毫秒级控制和日级研究计划意义不同，同样的候选数量在不同表示结构中也不等价。评估协议应保存任务族、时间尺度、容量基线与观测窗口，跨系统比较时报告置信区间，而不是直接比较一个抽象“阻力值”。只有相对基线稳定，区间转移才具有可复现含义。

\begin{table}[htbp]
\centering
\begin{tabularx}{\textwidth}{l X X X}
\toprule
\rowcolor{structurecolor!20} 区间 & 事实侧诊断 & 主体侧诊断 & 优先干预 \\
\midrule
顺行区 & 歧义低、步骤短、验证明确 & 目标一致、控制投入低 & 低开销执行并监测环境变化 \\
攻坚区 & 歧义或验证成本高 & 目标一致、愿意持续投入 & 补充证据、任务分解、外部工作台 \\
拖滞区 & 路径清楚、技术难度低 & 冲突或启动成本高 & 重构接口、自动化、承诺与卸载 \\
失稳区 & 不可辨识或搜索负担高 & 冲突、风险和负荷同时高 & 降低分辨率、补证、协作或安全退出 \\
\bottomrule
\end{tabularx}
\end{table}

\section{诊断闭环、干预边界与理论位置}
\label{sec:section_6127}

双重不对称性的目的不是制造一个总阻力比喻，而是把失败拆成可以干预的来源。诊断从事实侧开始：目标是否明确，观测是否足够，候选是否可辨识，约束是否相容，验证是否可执行；随后检查主体侧：候选是否合格，目标是否冲突，权限是否充足，资源是否超过预算，控制投入是否可持续。顺序并非绝对，但必须避免用主体偏好替代事实验证，也避免用事实难度否定目标重议的正当性。

我们把一次闭环记录为
\begin{equation}
\mathcal{T}_k=(x_k,n_k,S_k,\pi_k,a_k,o_{k+1},e_{k+1}),
\end{equation}
其中 $x_k$ 是事实状态估计，$n_k$ 是规范状态，$S_k$ 是候选集，$\pi_k$ 是策略，$a_k$ 是已提交行动，$o_{k+1}$ 是新观测，$e_{k+1}$ 是误差与外部反馈。该记录使我们能够判断失败发生在表示、搜索、资格筛选、执行还是反馈写回。只有写回真实结果，系统才有资格更新事实模型或主体策略；内部流畅度不能替代外部检验。

干预应对应诊断来源。不可辨识时增加传感、实验或来源独立的证据；搜索爆炸时改变分解、启发式或表示；接口丢失时恢复身份、版本和绑定；目标冲突时使用优先级、Pareto 比较或最小冲突集；资源过载时限流、卸载或延迟执行；授权不足时请求权限而非绕过边界。一次干预可能产生副作用，因此应记录前后任务成功、校准、资源、风险和外部影响，并保留回滚条件。

反例尤其重要。一个主体“很想做”但缺少关键证据，增加动机不会解决不可辨识；一个主体“会做”但拒绝越权，拒绝不是能力不足；一个系统快速给出答案，既可能因为任务简单，也可能因为跳过验证；一个系统停止执行，既可能是失稳，也可能是正确识别出硬约束冲突。双重模型要求这些情形在日志和测试中可区分，避免把所有行为压缩成单一强度轴。

在多主体环境中，两重不对称还会跨边界传播。一个 Agent 的事实证据可能成为另一个 Agent 的外部记忆，一个主体的规范状态却不能未经授权直接写入另一个主体。协作接口必须区分事实主张、证据来源、建议、请求、承诺和命令；接收方重新执行资格检查。由此，群体一致不证明事实正确，共同偏好也不抹去个体边界。HSF-HD 在这里提供一般状态动力学的分层语言，AgentOS、WorldOS 或群体系统可以用不同结构实现同一组接口契约。

本章与系统论、协同论、耗散理论、复杂系统理论和认知科学的关系，是继承开放系统、反馈、资源约束、涌现协作与有限认知等原则；这些原则不直接给出本章公式。HSF-HD 将其转化为事实候选收缩、规范资格筛选、联合调度和状态转移的模型族，再由 AgentOS 用 $h(t)$、$E$、$\Theta$、$\Psi$、SPC、TCE、CCM、Boundary、Offloading 与 $M_e$ 实现一个工程实例。这样的层级避免把一般理论和具体系统同义化。

最终，我们把“智力高山”和“欲望风口”保留为阅读直觉，却不让它们承担证明。正向与逆向的差异必须由观测核、候选结构和算法协议验证；顺行与逆行的差异必须由规范状态、资格集合和控制日志验证。智能系统真正要寻找的不是一条被神秘场扭曲的测地线，而是在事实证据充分、规范约束合法、资源投入可持续的条件下，构造能够执行、反馈、修订和解释的策略轨迹。

这种策略轨迹还必须允许明确失败：无法满足的约束应输出冲突证据，无法区分的候选应输出剩余歧义，资源不足应输出预算缺口。可解释的“不知道、不能做或暂不做”是模型边界的一部分，而不是智能状态动力学需要掩盖的异常。
